\documentclass[../../../../hsm_tok_q.tex]{subfiles}

\begin{document}

    \textsf{図 \ref{fig/hsm_tok_2025_2nd_01_09_01_q}}で，%
    $\mathrm{△ABC}$は、$\angle\mathrm{ACB}$が鈍角の三角形である．

    解答欄に示した図をもとにして，辺AB上にあり，
    辺ACと辺BCまでの距離が等しい点Pを，
    定規とコンパスを用いて作図によって求め，
    点Pの位置を示す文字Pも書け．
    
    ただし，作図に用いた線は消さないでおくこと．
    \\
    \begin{minipage}{\linewidth}
        \vspace{.5\baselineskip}
        {\centering
            \setlength{\abovecaptionskip}{1pt}
            \includegraphics%
                {\subfix{fig_hsm_tok_2025_2nd_01_09/fig_hsm_tok_2025_2nd_01_09_01_q.pdf}}
                \captionof{figure}{} \label{fig/hsm_tok_2025_2nd_01_09_01_q}
        }
    \end{minipage}

    \vspace{.2\baselineskip}
\end{document}

